% =============================================================================
% capitolo2_prototipo.tex
% Capitolo 2: Il prototipo - Astrazione e Funzionalita'
%
% LINEE GUIDA VITALI:
%  Questo capitolo e' dedicato al "COSA" e al "PERCHE'": NON al "COME".
%  Il lettore deve capire il sistema senza conoscere i dettagli tecnici.
%
%  A) PANORAMICA AD ALTO LIVELLO:
%     Cosa fa il sistema, a cosa serve, per chi e' pensato.
%     Un diagramma a blocchi (block diagram) e' molto utile qui.
%
%  B) CASI D'USO:
%     Scenari d'uso principali con attori e azioni.
%     Puoi usare tabelle, diagrammi UML use-case, o elenchi narrativi.
%     NON descrivere come sono implementati: solo cosa fanno.
%
%  C) INTERFACCIA UTENTE / API:
%     Mostra screenshot, wireframe, o esempi di API (se applicabile).
%     Spiega come si usa il sistema dal punto di vista dell'utente finale.
%
%  ATTENZIONE: Nessun dettaglio implementativo in questo capitolo.
%              Linguaggi, librerie, schemi DB vanno nel Cap. 3.
%  LUNGHEZZA SUGGERITA: 10-20 pagine.
% =============================================================================

% --- A. PANORAMICA ---
% Cosa fa il sistema ad alto livello, a cosa serve, per chi.
% Includi un diagramma a blocchi se possibile (con tikzpicture o immagine).

% --- B. CASI D'USO ---
% Elenca i principali use case: attore, precondizioni, flusso principale.
% NON spiegare COME sono implementati.

% --- C. INTERFACCIA / API ---
% Come si usa il sistema: screenshot, workflow utente, esempi di chiamate.

Scrivi qui il contenuto del capitolo 2 sul prototipo e le sue funzionalita'.
